\begin{pastexam}[label={px:gk3-111-ds-01}]{高考三級資訊處理 111 年 第一題}
（申論；25 分）

以下是一中序運算式（Infix expression）轉換（Convert）成後序運算式（Postfix expression）的演算法
\begin{minted}[fontsize=\small, frame=single]{text}
operstk = the empty stack;
while(not end of input){
  symb = next input character;
  if(symb is an operand)
       add symb to the postfix string;
  else{
    while(!empty(operstk) && precedence(stacktop(operstk),symb)){
      topsymb = pop(operstk);
      add topsymb to the postfix string;
       } /*end while*/
       if (empty(operstk) || symb != ')')
           push(operstk, symb);
        else
           topsymb = pop(operstk);
       } /*end else*/
   } /*end while*/
while(!empty(operstk)){
  topsymb = pop(operstk);
  add topsymb to the postfix string;
} /*end while*/
\end{minted}
其中資料結構：\\
“\texttt{operstk}”：用來儲存運算子的堆疊（Stack）；\\
“\texttt{stacktop(operstk)}”：表示 top 指標所指堆疊 \texttt{operstk} 的運算子；\\
程序（Procedures）或函數（Functions）：\\
“\texttt{empty(operstk)}”：檢查堆疊 \texttt{operstk} 是否為空的布林函數；\\
“\texttt{pop(operstk)}”：從堆疊 \texttt{operstk} 中取出一運算子；\\
“\texttt{push(operstk, symb)}”：將運算子 \texttt{symb} 存入堆疊 \texttt{operstk}；\\
“\texttt{precedence(op$_1$,op$_2$)}”：布林函數，定義在一沒有左右括弧的中序運算式中，op$_1$ 運算子
出現在 op$_2$ 運算子的左邊時，當 op$_1$ 運算子優先順序不低於 op$_2$ 運算子，則設定成 TRUE，否則為
FALSE。例如，我們給定 \texttt{precedence(\textquotesingle{}*\textquotesingle{}, \textquotesingle{}+\textquotesingle{})=TRUE}，\texttt{precedence(\textquotesingle{}+\textquotesingle{}, \textquotesingle{}+\textquotesingle{})=TRUE}，
\texttt{precedence(\textquotesingle{}+\textquotesingle{}, \textquotesingle{}*\textquotesingle{})=FALSE}，為了處理運算式左右括弧，設定下列的 precedence:
\begin{minted}[fontsize=\small]{text}
precedence('(', op) = FALSE /*op 為任一運算子*/
precedence(op, '(') = FALSE /*op 為除')'外的任一運算子*/
precedence(op, ')') = TRUE  /*op 為除'('外的任一運算子*/
precedence(')', op) = undefined   /*op 為任一運算子*/
\end{minted}
以中序運算式(2+3)*4 為例，執行上述演算法，依處理每一個運算子或運算元時，輸出 postfix string 及
operstk 內容為何（“eos”表示 end of string）？（25 分）
\begin{center}
\begin{tabular}{ccc}
\toprule
symbol & \qquad postfix string \qquad & \qquad operstk \qquad \\
\midrule
( & & \\
2 & & \\
+ & & \\
3 & & \\
) & & \\
* & & \\
4 & & \\
eos & & \\
\bottomrule
\end{tabular}
\end{center}
\end{pastexam}

\begin{pastsoln}
原卷的程式和 precedence 定義用彎引號把字元括起來，這裡改成 C 語言的直引號，意思一樣。

\textbf{規則整理。} 這個演算法就是資結書虛擬碼~6.8 的中序轉後序，只是把「什麼時候 pop」寫成
\texttt{precedence(stacktop, symb)}：
\begin{itemize}
\item 運算元直接接到輸出。
\item 讀到運算子或括號 \texttt{symb} 時，只要堆疊頂端的運算子「優先順序不低於」\texttt{symb}，
就 pop 出來接到輸出。 左括號在堆疊裡時 \texttt{precedence(\textquotesingle{}(\textquotesingle{}, op)} 是 FALSE，所以左括號像一道牆，
不會被一般運算子 pop 掉；讀到 \texttt{\textquotesingle{})\textquotesingle{}} 時 \texttt{precedence(op, \textquotesingle{})\textquotesingle{})} 對 \texttt{\textquotesingle{}(\textquotesingle{}} 以外的
運算子都是 TRUE，會一路 pop 到遇到 \texttt{\textquotesingle{}(\textquotesingle{}} 為止。
\item 接著若堆疊是空的或 \texttt{symb} 不是 \texttt{\textquotesingle{})\textquotesingle{}}，就把 \texttt{symb} push 進去；否則
（\texttt{symb} 是 \texttt{\textquotesingle{})\textquotesingle{}}，頂端一定是對應的 \texttt{\textquotesingle{}(\textquotesingle{}}）把 \texttt{\textquotesingle{}(\textquotesingle{}} pop 掉丟棄，
括號不會出現在輸出裡。
\item 讀完輸入（eos），把堆疊剩下的運算子依序 pop 到輸出。
\end{itemize}

\textbf{追蹤。} operstk 由底寫到頂（最右邊是頂端）。
\begin{center}
\begin{tabular}{cll>{\raggedright\arraybackslash}p{6.2cm}}
\toprule
symbol & postfix string & operstk & 說明 \\
\midrule
( & （空） & ( & 堆疊空，直接 push \\
2 & 2 & ( & 運算元，輸出 \\
+ & 2 & ( + & \texttt{precedence(\textquotesingle{}(\textquotesingle{}, \textquotesingle{}+\textquotesingle{})} 為 FALSE，不 pop；push \texttt{+} \\
3 & 23 & ( + & 運算元，輸出 \\
) & 23+ & （空） & \texttt{precedence(\textquotesingle{}+\textquotesingle{}, \textquotesingle{})\textquotesingle{})} 為 TRUE，pop 出 \texttt{+}；頂端是 \texttt{\textquotesingle{}(\textquotesingle{}}，
\texttt{precedence(\textquotesingle{}(\textquotesingle{}, \textquotesingle{})\textquotesingle{})} 為 FALSE，停；\texttt{symb} 是 \texttt{\textquotesingle{})\textquotesingle{}}，pop 掉 \texttt{\textquotesingle{}(\textquotesingle{}} 丟棄 \\
* & 23+ & * & 堆疊空，push \\
4 & 23+4 & * & 運算元，輸出 \\
eos & 23+4* & （空） & 輸入結束，pop 出 \texttt{*} \\
\bottomrule
\end{tabular}
\end{center}
最後的後序運算式是 \texttt{23+4*}，也就是 $(2+3)\times 4 = 20$。 檢查：後序求值時先算
$2\ 3\ +$ 得 5，再算 $5\ 4\ *$ 得 20，和原式一致。
\end{pastsoln}
